\section{CONCLUSION}\label{conclusion}

Across a 24-LLM-configuration evaluation matrix on n=125
Synthea-generated FHIR Conditions, fabrication of non-existent
ICD-10-CM codes ranged from 0\% to 5.6\% for flagship models (Opus
4.7, GPT-5.5, Gemini 2.5 Pro) and from 0\% to 12.8\% for sub-flagship
models. Constrained prompting eliminated fabrication entirely for
every Anthropic and OpenAI model. Hallucination behaves as a
controllable property of the prompting setup at this scale rather
than a fixed model limitation. Abstention, separated from
fabrication via the 6-way taxonomy, proved the more consequential
failure mode for some models.

Two caveats temper the headline. The n=125 cohort and the narrow
12-code ACCESS-Model vocabulary limit per-cell precision and bound
the population that surfaces fabrication. The D4
abbreviation-substitution stress is only one proxy for real-world
hallucination triggers; multi-condition narratives, ambiguous
documentation, code-specificity ambiguity, and long-tail rare
diagnoses are not exercised here. A v2 cohort with broader
vocabulary and more realistic clinical text (e.g., discharge
summaries, ED notes) is needed to establish the true effect size
of fabrication risk on production-relevant inputs.

Two deployment prescriptions emerge. \textbf{Zero-shot LLM prompting is
unsafe for medical concept normalization}, since every
fabrication observed occurred under zero-shot, and constrained or
retrieval-augmented grounding is required to eliminate it. Among
grounded configurations, \textbf{cost-per-correct rather than accuracy
alone should drive selection}, because clinical coding errors are
downstream-costly and the cheaper model can win on the combined
metric. The deployment frontier consists of \textbf{GPT-4o-mini
constrained} (94.2\% top-1, \$0.0003 per correct, the cheapest
configuration at $\geq$94\% accuracy) and \textbf{Claude Haiku 4.5
constrained} (96.8\%, \$0.0044 per correct, an additional 2.6pp
accuracy at roughly 14× the cost). Both dominate the largest
frontier model (Opus 4.7 constrained, 94.8\%, \$0.0133). The
supported workflow is LLM-augmented coding with terminologist
review.
